\section{Application}\label{sec:application}

The COVID-STEROID 2 trial was a multi-center study comparing daily dose 6 mg (low) of dexamethasone for up to 10 days vs.\ 12 mg (high) in patients with severe and critical COVID-19. The full description of the study including all results can be found in \citep{10.1001/jama.2021.18295}. The primary outcome of the trial was days alive and without out use of life-support in the period from randomization to day 28. In other words, that outcome is an integer between 0 and 28 with a substantial peak at day 0 corresponding to the atom $\mathcal{E}$ coming from the patients who are never taken off life-support systems. The analyses in the main publications \citep{10.1001/jama.2021.18295} were conducted using the test proposed in this paper. In the JAMA publication the test is being referred to as the Kryger Jensen and Lange test. In this section we re-analyze the sub-group consisting of patients in need of invasive mechanical ventilation at baseline ($n = 206$). These can generally be taken to be the most severely ill patients.

The distribution of the outcome in each randomization group is presented by histograms in Figure~\ref{fig:application} (left and middle panels). In our notation the outcome (days alive without life support) is re-coded such that $A$ takes the value zero if the outcome is zero or one otherwise. $Y$ is said to be the observed number of days without life support.

\begin{figure}[htbp]
\centering
\includegraphics[width=\textwidth]{application.pdf}
\caption{Histograms of days alive without life support from the COVID-STEROID 2 trial stratified by treatment dose (left, middle) and simultaneous confidence regions based on the semi-parametric likelihood ratio test (right).}
\label{fig:application}
\end{figure}

Comparing the groups by a Wilcoxon test yields a p-value of \ApplicationWilcoxP{} hence leading to a non-significant conclusion at the pre-specified 5\% level. Using our proposed parametric likelihood ratio test yields a p-value of \ApplicationLRTP{} and therefore leads to a rejection of the null hypothesis. From the histograms it could be questioned whether the continuous part of the distributions are normally distributed. Accordingly, we also apply our proposed semi-parametric version of the likelihood ratio test. This yields a p-value of \ApplicationSPLRTP{}. In conclusion, our proposed method is able to detect the difference between the groups. Further and in contrast to the Wilcoxon test, our method also provides a way to interpret the difference between the groups. This is illustrated in the right panel of Figure~\ref{fig:application} showing simultaneous confidence regions for the two effect parameters. In this case it is clear that the primary effect of the intervention is not on the discrete component (i.e., the proportion of patients who never get off life support). Instead, the effect appears to be on the continuous part and is between half a day and seven days.
